\chapterimage{figs/chapterhead.png}
\chapter{Example models and initial use of the simulator}
\label{chap:modelos_de_exemplo_e_uso_inicial}

\section{Introduction}
\label{sec:cap4_introducao}

Once the \textit{GenESyS} graphical interface is up and running and the reader has an initial view of your organization, the next natural step is to start working with concrete models. This is the point at which the simulator stops being just an application open on the screen and starts to become an effective modeling and experimentation environment. To achieve this, the most appropriate path is not to start immediately by creating a complete model from scratch, but to use the example models already available in the project.

This decision is important from both a didactic and practical point of view. From a didactic point of view, an example model already offers a coherent structure, with interconnected elements, completed properties and expected behavior. This reduces the initial cognitive load on the user, who does not need to simultaneously learn the interface, the semantics of the components, the flow construction logic and the simulation execution mechanism. From a practical point of view, example models allow you to check early on whether the system installation is functional and whether the GUI can open, display, interpret and execute real models.

In this chapter, the objective is to teach the reader how to use the models saved in the \texttt{models/} folder as a gateway to the simulator. At the end of reading, the user should know how to locate an example model, load it into the graphical interface, read its visual structure, identify the main flow of the modeled system, run the simulation and make small controlled modifications to observe the effect of these changes.

\section{Why start with example templates}
\label{sec:cap4_por_que_comecar_por_modelos_exemplo}

Rich simulations tend to offer many job possibilities from the first contact. This wealth, although valuable, can become an obstacle when a good entry strategy is not chosen. In the case of \textit{GenESyS}, using example templates solves exactly this problem.

When opening an already saved model, the reader has immediate access to an artifact that:
\begin{itemize}
\item is already organized coherently;
\item already contains a flow logic or modeling structure;
\item can now be inspected by the GUI;
\item can now be executed;
\item can now be used as a basis for tracked changes.
\end{itemize}

This allows learning to occur in layers. First, the user learns to recognize the model as a visual object within the interface. Then learn how to execute it. Then, begin relating the graphic elements to their effects in the simulation. Only later did it begin to produce its own models with greater autonomy.

In pedagogical terms, the example model fulfills a role similar to that of an experiment already set up in a laboratory: it does not replace active learning, but creates a stable initial context in which the user can begin to observe, formulate hypotheses, modify parameters and interpret results.

\subsection{The mistake of starting from scratch too soon}
\label{subsec:cap4_erro_comecar_do_zero}

There is a natural temptation, especially in technically curious readers, to try to build their own model during their first interactions with the system. Although this initiative may seem productive, it often introduces unnecessary difficulties.

When the user starts too early from an empty space, he needs to decide at the same time:
\begin{itemize}
\item which elements to insert;
\item in what order to connect them;
\item which properties to adjust;
\item what supporting data is needed;
\item how to check if the model is coherent;
\item what to expect from execution.
\end{itemize}

This accumulation of decisions makes it more difficult to distinguish what is a difficulty inherent to the modeling and what is simply an initial lack of knowledge of the tool. The example model avoids this noise. It separates the problem of using the interface from the problem of full model authoring.

\subsection{Example templates as reading material}
\label{subsec:cap4_modelos_exemplo_como_leitura}

In addition to serving as executable cases, the example templates should also be read. This reading, however, is not the same as reading a text or a source code. It involves:
\begin{itemize}
\item observe the visual arrangement of the model;
\item perceive the main flow sequence;
\item identify entry and exit points;
\item recognize auxiliary elements that support execution;
\item report visual structure to observed behavior.
\end{itemize}

With this, the model stops being a file and becomes an object of study within the GUI.

\section{The \texttt{models/} folder as a reference for use}
\label{sec:cap4_pasta_models_referencia}

For the \textit{GenESyS} user, the \texttt{models/} folder should be understood as one of the most important areas of the project. This is where the saved models that will be used as an entry point for exploring the simulator are located. In this first part of the book, this folder has much higher priority than any source code directory.

The role of this folder is twofold. First, it concentrates concrete material for using the simulator. Second, it provides the user with a repertoire of examples from which to learn how to open, read, run, and change models. This means that for initial learning purposes, the user does not need to worry about software implementation details. Your main universe of work is precisely in these model files.

\subsection{What to look for in the \texttt{models/} folder}
\label{subsec:cap4_o_que_procurar_models}

When browsing the \texttt{models/} folder, the user must look for models that preferably satisfy three criteria:

\begin{enumerate}
\item are visually simple or moderately simple;
\item have an easily identifiable main flow;
\item allow an execution whose behavior can be observed without great interpretative effort.
\end{enumerate}

This doesn't mean that larger or fancier models aren't valuable. It just means that, when entering the simulator, it is advisable to prioritize cases that favor clarity and readability.

\subsection{How to select good first examples}
\label{subsec:cap4_como_selecionar_bons_exemplos}

As a first approximation, the best examples tend to be those where the user can answer the following questions relatively quickly:
\begin{itemize}
\item where does the flow begin?
\item what main elements does it go through?
\item where does it end?
\item Are there queues, resources, attributes, or auxiliary data that support this behavior?
\item what do I expect to see when this model is run?
\end{itemize}

If these questions can be answered without great difficulty, the model is probably a good candidate for initial exploration.

\section{Opening a sample model in the GUI}
\label{sec:cap4_abrindo_modelo_exemplo}

With the application already started, the next step is to open a real model. The operation of opening a model file should not be treated as a simple technical detail, because it inaugurates the phase in which the GUI starts to operate effectively on a concrete work object.

The general procedure is as follows:
\begin{enumerate}
\item start the \textit{GenESyS} GUI;
\item find the menu or command corresponding to opening the model;
\item navigate to the \texttt{models/} folder;
\item select an appropriate template file;
\item confirm the opening and observe the interface update.
\end{enumerate}

After that, the central graphics area should start displaying the model, and the GUI panels will tend to reflect the loaded elements, allowing for more detailed inspection.

\subsection{What to confirm immediately after opening}
\label{subsec:cap4_confirmar_apos_abertura}

Once the template is loaded, the user must confirm a few simple but important points:

\begin{itemize}
\item the graphics area has been filled with the model structure;
\item there are visible elements and connections;
\item the interface seems to recognize the model as open;
\item inspection and simulation commands became usable;
\item the user can visually navigate the loaded content.
\end{itemize}

This confirmation is useful because it distinguishes two very different states of the application:
\begin{itemize}
\item the GUI is empty or without an active model;
\item the GUI effectively working on a model.
\end{itemize}

Only in the second case does it make sense to proceed to reading and execution.

\subsection{When the model appears ``too big''}
\label{subsec:cap4_modelo_grande_demais}

In some cases, the first model opened may appear visually too dense. If this occurs, the user should not immediately interpret this sensation as a problem with the simulator or model. In general, this just indicates that:
\begin{itemize}
\item zoom needs to be adjusted;
\item the chosen model was not the most suitable for the first exploration;
\item reading must start with a main region, and not with the entire scene.
\end{itemize}

In this situation, the best practice is to return to a simpler example or restrict the analysis to the most easily identifiable main flow.

\section{How to read a model in the graphics area}
\label{sec:cap4_como_ler_modelo_area_grafica}

Once the model is loaded, the user's next task is to read it. This reading is neither purely visual nor purely technical. It is a reading guided by questions about behavior and structure.

The best initial reading order is:
\begin{enumerate}
\item find the beginning of the flow;
\item follow the main path of the model;
\item identify the end of the flow;
\item observe support elements;
\item only then examine local properties and details.
\end{enumerate}

This order is important because it prevents the user from getting lost in details before understanding the central logic of the model.

\subsection{Identifying the main flow}
\label{subsec:cap4_identificando_fluxo_principal}

The main flow is the most important structure of the model for first contact. It corresponds to the essential path taken by the main entities, events or transformations represented in the system. In many cases, this flow can be recognized by a sequence of components connected in a relatively linear way or by a visual arrangement that clearly suggests the progression of the modeled process.

The question that should guide the user is:
\begin{quote}
What is the minimal story this model is telling?
\end{quote}

When this question begins to be answered, the model stops being a set of elements on the screen and starts to be perceived as an organized system.

\subsection{Distinguishing main flow from auxiliary elements}
\label{subsec:cap4_fluxo_principal_elementos_auxiliares}

A careful reading of the model depends on distinguishing what constitutes its central path from what supports it. Often, in addition to the main flow, the model also involves auxiliary data and structures: queues, resources, statistics, attributes, variables, counters or other definitions.

The user should not ignore these elements, but he should not start with them either. The correct procedure is:
\begin{enumerate}
\item understand the flow first;
\item then ask what support elements allow this flow to work;
\item only then observe specific properties.
\end{enumerate}

This order greatly improves the readability of the model.

\section{Running the model for the first time}
\label{sec:cap4_executando_modelo_primeira_vez}

After opening and reading the model, the decisive moment comes: running the simulation. It is at this stage that the user checks whether the observed visual structure actually corresponds to understandable dynamic behavior.

On first use, execution should be conducted with exploratory intent and not in a rush. The objective is not just to ``run'' the model, but to understand:
\begin{itemize}
\item if the simulation starts correctly;
\item whether the model's behavior makes sense in relation to its structure;
\item whether the GUI provides sufficient signals to track progress;
\item whether the model reaches a final state or evolves coherently.
\end{itemize}

\subsection{What to watch out for while running}
\label{subsec:cap4_o_que_observar_durante_execucao}

When running a model for the first time, the user should avoid trying to read all output and all panels at the same time. It's best to focus on a few core aspects:

\begin{itemize}
\item there is a clear indication that the simulation has started;
\item the model appears to visually evolve or produce coherent signs of progress;
\item the observed behavior is compatible with the previously identified main flow;
\item there is no obvious structural flaw in the execution process.
\end{itemize}

In other words, the first run serves to connect the static model to its dynamics.

\subsection{First validation of behavior}
\label{subsec:cap4_primeira_validacao}

A well-done initial validation is modest but sufficient. The user must be able to respond:
\begin{itemize}
\item did the model open correctly?
\item Could the simulation be started?
\item does the observed behavior seem plausible?
\item Did the execution finish or progress without obvious inconsistencies?
\end{itemize}

If these answers are positive, then the first practical contact with the simulator has already been successful.

\section{Modifying an example model in a controlled way}
\label{sec:cap4_modificando_modelo_controlado}

After opening and running a model, the next most productive step is to make small, controlled modifications. This step is pedagogically very important, because it transforms the example model into an active learning instrument.

The best initial modifications are those that:
\begin{itemize}
\item do not completely change the structure of the model;
\item can be easily reversed;
\item produce an observable effect;
\item help the user understand the function of properties or parameters.
\end{itemize}

Examples of appropriate changes at this stage include:
\begin{itemize}
\item adjust a simple property;
\item change a name or label;
\item change a numerical value associated with a component;
\item save a modified copy of the template.
\end{itemize}

\subsection{Why small changes are better}
\label{subsec:cap4_por_que_pequenas_mudancas}

If the first modification is too big, the user loses the ability to report cause and effect. It is difficult to know whether the observed change in behavior is due to a specific parameter, an extensive structural change or an inconsistency introduced accidentally.

Small changes, on the contrary, allow the user to experience the simulator logic almost like in a laboratory:
\begin{enumerate}
\item observe the original state;
\item change a specific point;
\item rerun;
\item compare the effect.
\end{enumerate}

This methodology is much more useful than making extensive modifications at the beginning.

\section{Best practices for exploring example models}
\label{sec:cap4_boas_praticas_explorar_modelos}

A few simple practices make initial use of the example templates much more beneficial.

\subsection{Start with smaller cases}
\label{subsec:cap4_comece_por_casos_menores}

Smaller models generally allow for clearer reading of the interface, flow, and behavior. They also facilitate the association between visual structure and observed execution.

\subsection{Save separate copies}
\label{subsec:cap4_salve_copias_separadas}

When modifying a sample template, it is recommended that you save a new copy rather than immediately overwriting the original file. This preserves the reference example and makes later comparisons easier.

\subsection{Use execution as a reading tool}
\label{subsec:cap4_execucao_como_leitura}

Often, the model only becomes fully intelligible when it is executed. Therefore, execution should not be seen just as a final step, but as part of the process of understanding the example itself.

\subsection{Return to structure after running}
\label{subsec:cap4_volte_a_estrutura_apos_executar}

A very useful practice is to re-read the model visually after the first run. Observed behavior often sheds new light on the function of its elements, connections, and supporting data.

\begin{table}[htbp]
\centering
\caption{Recommended workflow for exploring example models in GenESyS.}
 \label{tab:cap4_roteiro_modelos_exemplo}
\small
 \begin{tabular}{|p{4.2cm}|p{7.1cm}|}
\hline
\textbf{Step} & \textbf{Goal} \\
\hline
Open the model & Load a working case in the GUI \\
\hline
Read the main flow & Understand the central logic of the modeled system \\
\hline
Recognize supporting elements & Understand what data and structures support the flow \\
\hline
Run the simulation & Report visual structure to observed behavior \\
\hline
Change something simple & Turn the example into an active learning object \\
\hline
Rerun and compare & Observe the effect of the modification made \\
\hline
 \end{tabular}
\end{table}

\section{Final Considerations}
\label{sec:cap4_consideracoes_finais}

In this chapter, example models have been presented as the user's main entry point into practical work with \textit{GenESyS}. Instead of starting with the creation complete of a new model, the reader was instructed to open already saved cases, read them visually, identify their main flows, execute them and make small controlled changes. This strategy allows learning to happen in a progressive and concrete way, without requiring the user, right from the beginning, to fully master the entire tool.

The essential point to remember is that an example model should not be treated as a simple passive demonstration. It is, at the same time, an object of reading, an object of execution and an object of experimentation. From here, the user is not only familiar with the GUI as a structure, but begins to use it effectively on real models. The next natural step is to deepen the relationship between graphical interface, execution, behavior observation and textual representation of the model, which will be done in the next chapter by introducing the \textit{Simulang} tab, the execution observation mechanisms and the first procedures for analyzing the system's behavior.

Test your understanding of this content by answering the following questions:

\begin{enumerate}
\item Why are example templates a better entry point than immediately creating a complete template from scratch?

\item What is the most productive order to read a model loaded in the GUI?

\item What should the user validate when first running a sample model?

\item Why is it important to initially modify only small aspects of the model?
\end{enumerate}
